\begin{table}[htbp]
\centering
\caption{Who bears a fiscal gap}\label{tab:whopays}
\footnotesize
\begin{tabularx}{\textwidth}{>{\raggedright\arraybackslash}X>{\raggedright\arraybackslash}X>{\raggedright\arraybackslash}X>{\raggedright\arraybackslash}X}
\toprule
Route & Who pays & Decided by & Price or size in this paper \\
\midrule
Fiscal adjustment & Taxpayers or program beneficiaries & Legislation & Threshold $\varphi^*$: about 0.47 of marginal interest cost (United States, 2023–25) \\
Inflation, tolerated by the central bank & Holders of nominal claims: savers, bondholders, currency holders (part of them abroad) & Central bank, without a legislative vote & About 0.2 pp of inflation a year for a decade per 0.1 of missing offset; Japan 2022–25: about 18\% of GDP from holders of nominal claims \\
Higher real rates, when the central bank tightens & Taxpayers, through interest paid to bondholders and, on reserves, to banks; workers, through slower growth & Central bank & Above $\kappa^*$ = 1/R (0.46 in the United States in 2025), inflation cannot close the gap at all; United States and UK 2021–25: most of a 10–14\% of GDP transfer returned; UK interest on reserves alone 5.7\% of GDP \\
\bottomrule
\end{tabularx}
\end{table}
